% 胡不归模型主图：直线 BC, 定点 A 在一侧；P 在 BC 上；在 B 作角 α 使 sinα=k，过 P 作 PH⊥l'
\documentclass[tikz,border=4pt]{standalone}
\usepackage{ctex}
\usepackage{amsmath}
\usetikzlibrary{calc, arrows.meta, angles, quotes}
\begin{document}
\begin{tikzpicture}[scale=1.0, thick]
  \coordinate (B) at (0, 0);
  \coordinate (C) at (7, 0);
  \coordinate (A) at (1.8, 3.0);
  % 取 α=30°，l' 在 BC 下方与 BC 成 30°
  % l' 方向：(cos(-30°), sin(-30°)) = (0.866, -0.5)
  \coordinate (Ll) at (6.5, -3.75);  % 沿 l' 远端
  % P 取在 BC 上某位置（中间）；从 P 向 l' 作垂线
  \coordinate (P) at (3.6, 0);
  % P 到 l' 的垂足 H：l' 方向 (0.866, -0.5)，l' 通过 B(0,0)
  % H = P 在 l' 上的投影：t = P·dir = 3.6*0.866 = 3.118
  % H = (3.118 * 0.866, 3.118 * (-0.5)) = (2.700, -1.559)
  \coordinate (H) at (2.700, -1.559);

  \draw[thick] (-0.5, 0) -- (7.5, 0);
  \node[right] at (7.5, 0) {$BC$};
  \draw[thick, blue] (B) -- (Ll);
  \node[blue, right] at (6.5, -3.75) {$l'$};

  \filldraw (A) circle (2pt) node[above] {$A$};
  \filldraw (B) circle (2pt) node[above left] {$B$};
  \filldraw (P) circle (2pt) node[above] {$P$};
  \filldraw (H) circle (2pt) node[below right] {$H$};

  \draw[red, thick] (A) -- (P);
  \draw[red, thick] (P) -- (H);

  % 角 α
  \draw[blue, thin] (1.2, 0) arc[start angle=0, end angle=-30, radius=1.2];
  \node[blue, font=\footnotesize] at (1.45, -0.35) {$\alpha$};

  % 直角符号（PH⊥l'）
  \draw[blue, thin]
    ($(H) + (-0.15,-0.26)$) -- ($(H) + (0.10, -0.43)$)
    -- ($(H) + (0.26, -0.17)$);

  \node[red, font=\footnotesize] at (2.5, 1.7) {$PA$};
  \node[red, font=\footnotesize] at (3.5, -0.85) {$PH$};
  \node[font=\footnotesize] at (4.7, -1.7) {\small $PH = PB\sin\alpha = k\cdot PB$};
\end{tikzpicture}
\end{document}
